\chapter{Introduction}
\label{ch:intro}

\section{Motivation}
\label{sec:motivation}

A manipulator teleoperated through a single grounded haptic interface --- a device whose handle is mechanically fixed to the desk rather than carried freely in the hand, the very property that lets it render forces back onto the operator --- is used precisely where full autonomy is not yet trusted, or not yet possible, and where a remote operator's judgement must still resolve which of several objects to reach for, which side of an obstacle to pass, and when a task-level decision is required that no perception pipeline is asked to make~\cite{jain2020probabilistic}. Bilateral teleoperation architectures close this loop in both directions, forwarding the operator's motion to the manipulator and rendering the resulting interaction back to the operator's hand~\cite{hokayem2006bilateral}. The interface itself imposes a structural asymmetry, established by Kim et al.'s comparison of the two dominant motion mappings~\cite{kim1987comparison}: its own workspace is a fraction of the manipulator's reach, so a position-controlled mapping requires the operator to periodically disengage and reposition the handle, while a rate-controlled mapping removes that limitation at the cost of a less intuitive, integrating relationship between hand and end-effector. Extending such a single-handed interface to a bimanual manipulator multiplies both the coordination and the risk: the operator drives one arm at a time through the one channel available, while both arms must remain clear of the environment and of each other throughout.

Shared autonomy relieves part of this load by inferring, from the operator's own command, which of several candidate goals is intended, and by combining an autonomous policy toward that goal with the operator's raw input~\cite{selvaggio2021autonomy,losey2018review}. Safety, in the sense used throughout this thesis, means that the motion actually commanded to the arms satisfies hard constraints at every control tick --- no collision with the modelled obstacles, including the two arms against each other and against themselves, and joint position and velocity limits respected. The canonical realization of shared autonomy is a convex blend of the two commands, weighted by an estimate of confidence in the inferred goal~\cite{dragan2013policy}. Such a blend does not, by itself, respect any constraint that the operator's command and the autonomous policy each individually satisfy: once the admissible set is nonconvex, as it generically is around real obstacle geometry, a blend of two individually safe commands can violate a constraint that neither command violates on its own, shown formally by Chaubey et al. and empirically by Guler et al.~\cite{chaubey2026minimalintervention,guler2026barrierik}. Confidently wrong assistance is worse than none: a blend that is safe only in expectation, and not enforced to be safe at the instant it is executed, can still drive an arm into contact.

A hard constraint applied after arbitration, rather than a soft penalty folded into the blend objective, is the response this literature converges on, most often realized as a control barrier function (CBF) embedded in a quadratic program (QP) solved at every control tick~\cite{ames2017cbf}. Published realizations of this post-arbitration principle remain partial: the closest, Guler et al.'s filtered pipeline, is validated only in simulation and a virtual-reality study, inside a position-space inverse-kinematics solve rather than a joint-velocity QP; the closest use of the filter's own dual variables as operator feedback, Raphalen et al.'s, is demonstrated around a single arm and without arbitration against an inferred autonomous policy~\cite{raphalen2026feeling}. No published system realizes a post-arbitration hard-safety filter, natively formulated at the joint-velocity level, on a physical bimanual manipulator, nor uses that filter's own shadow prices for anything beyond a one-way feedback channel.

This thesis addresses that combination directly: a strictly convex control-Lyapunov-function (CLF)--CBF QP enforces collision avoidance and joint limits as hard constraints of a nominal barrier condition on the command that finally reaches each arm, after an alignment-based arbitration law has already combined the operator's instantaneous twist with a belief-weighted autonomous policy. The controller and the perception pipeline supplying its collision geometry are validated end to end on a physical bimanual platform; the factorial comparison across teleoperation mappings and assistance strategies is then run at human-subject-study scale in a matched simulation environment.

\section{Research questions and hypotheses}
\label{sec:rqs}

The motivation above resolves into six research questions.

\begin{enumerate}[nosep, label=\textbf{RQ\arabic*.}]
\item (safety) Does the post-arbitration filter enforce its nominal barrier condition on physical hardware, and where does the sampled, physical closed loop depart from that condition?
\item (assistance channel) Which assistance channel, force guidance or reference blending, drives objective performance and subjective experience, and does combining both change either?
\item (motion mapping) Does the choice of teleoperation mapping interact with the assistance channel on objective and subjective measures?
\item (arbitration) Does alignment-based arbitration preserve operator primacy while still assisting?
\item (adaptive scheduling) Does reusing the QP's own shadow price to schedule tracking slack behave as intended when exercised on hardware?
\item (grasp-object barrier management) Does the hybrid grasp lifecycle move a target object into and out of the collision model without a detectable safety cost?
\end{enumerate}

A human-subject study tests RQ2 and RQ3 against four hypotheses, formed before any trial is run and stated here on named measures so that each is falsifiable against a quantity rather than an unstated notion of merit.

\begin{enumerate}[nosep, label=\textbf{H\arabic*.}]
\item (channel--mapping interaction) H1 predicts an interaction rather than one channel dominating: force guidance outperforms reference blending on task time $T$~\eqref{eq:mT} under the clutch mapping, and reference blending outperforms force guidance under the joystick mapping. The clutch-side half rests on a system-specific conjecture --- no prior source evaluates haptic guidance as a substitute for depth perception --- that a force pulling the handle toward the goal, with the handle visibly following it, compensates for the depth cues the rear-facing view denies the operator. The joystick-side half follows a channel-congruence argument instead: a guidance force adds to the effort of holding a direction or speed against the mapping's own centring spring, whereas reference blending arbitrates directly on the same velocity command the joystick already issues, so the two channels are expected to compete for authority under one mapping and to reinforce each other under the other, consistent with Abbink, Mulder, and Boer's principle that assistance should bias, never replace, the operator's own input~\cite{abbink2012hapticshared}.
\item (mapping) H2 predicts that the clutch mapping yields the shorter task time $T$ and path length $L$~\eqref{eq:mL}, and the same direction on the subjective composite index $\bar{S}$: the one-to-one correspondence between hand motion and gripper motion is expected to outweigh the cost of the re-indexing it requires --- itself guided in time by the joint-limit vibrotactile cue and in orientation by the alignment torque~\eqref{eq:aligntorque} --- against the joystick mapping's indirect, rate-controlled correspondence.
\item (authority cost of blending) H3 predicts that reference blending is rated \emph{worse} than force guidance alone on the workload subscales for mental and physical demand, $S_{\mathrm{md}}$ and $S_{\mathrm{pd}}$: a reference that no longer moves exactly as commanded registers as effort spent resisting or correcting it, an authority cost of the kind Zhang, Tron, and Khurshid identify wherever a felt intervention alters an operator's command rather than merely reporting it~\cite{zhang2021hapticagreement}. This effect is expected to be sharpest in the more cluttered relay scenario, where the belief-weighted policy's own safety margin can pull the blended reference away from an operator command that holds alignment on one task component, orientation for instance, while intending to pass closer to an obstacle than the policy allows; this scenario-specific half is noted for completeness, since the design averages the two scenarios into one value per condition and cannot test it.
\item (quality over speed) H4 predicts that, because a blended reference is one the safety filter can execute rather than one it must partially refuse, blending lowers path length $L$, mean tracking slack $\bar{\delta}$~\eqref{eq:mdelta}, and filter active time $T_{\lambda}$~\eqref{eq:mTlam} relative to force guidance alone, with no predicted direction on task time $T$, since nothing in the mechanism ties the size of the correction to the speed at which the operator reaches the goal.
\end{enumerate}

RQ4 through RQ6 are addressed by implementation and by exercise on physical hardware, not by a controlled experiment, and carry no tested hypothesis.

\section{Contributions}
\label{sec:contributions}

\begin{itemize}[nosep]
\item \textbf{A shared-autonomy framework combining autonomous grasping, a post-arbitration safety filter, haptic feedback, and command blending.} A strictly convex CLF--CBF QP, solved at the velocity-level inverse-kinematics (IK) stage with a CLF task-tracking constraint and a CBF collision constraint aggregated over the closest obstacle pairs by a soft minimum, enforces collision avoidance and joint limits on the already-arbitrated command reaching each arm. Authority given to the autonomous policy is arbitrated on the instantaneous directional agreement between the operator's commanded twist and the goal-directed policy's own twist, computed from a discounted belief score over candidate goals rather than a scalar confidence, distance, or entropy summary, and ramps to zero exactly when the two conflict, so the operator retains sole authority over any motion it opposes. The safety filter's own dual variables, already computed at every solve, are used as the robot's self-assessment of how constrained its surroundings are, and are low-pass filtered and fed back into the QP's own cost to schedule its per-arm tracking slack: tracking tightens in free space and yields near an active constraint, so precision and execution time are given up only where safety requires it, closing this loop internally rather than only rendering the multiplier to the operator as a feedback force. The same signal also drives, in development runs outside the reported evaluation, the time scaling of assigned trajectories, and opens the scheduling of the CLF convergence rate and of the posture weight. A grasped object is treated as a hybrid, not a continuously invariant, barrier: the target leaves and rejoins the declared obstacle pair set through explicit regimes --- exclusion during the blind approach, relaxed-margin contact during closure, re-parenting into the protected chain, and ramped reactivation on release --- each with its own reset map and switching condition, with continuity claimed only for the margin parameter within a regime, not for invariance of the original obstacle set across regimes.
\item \textbf{Deployment of the framework on physical bimanual hardware and in a matched simulation environment.} The controller and the camera-based perception pipeline supplying its collision geometry are validated end to end on a physical bimanual mobile manipulator, from camera-perceived obstacle geometry through teleoperated grasping under barrier-enforced collision avoidance. The same architecture also runs in a matched simulation environment built to exercise the framework at human-subject-study scale.
\item \textbf{A factorial human-subject study for validation.} The evaluation crosses both classical teleoperation mappings, position control with indexing and spring-centred rate control, with the assistance strategies above, over two scenarios that differ in whether they force the two arms to cooperate, run by human participants to test the framework's objective and subjective effect.
\end{itemize}

\section{Thesis outline}
\label{sec:outline}

The remainder of this thesis is organized as follows. Chapter~\ref{ch:sota} reviews the four literatures this thesis draws on --- teleoperation of remote manipulators, shared control and shared autonomy, safety-critical control, and perception for safety-critical control --- and closes by positioning the contributions above against the closest prior systems (Section~\ref{sec:sota-positioning}). Chapter~\ref{ch:method} develops the methodology: the QP-CLF-CBF safety controller (Section~\ref{sec:qp}), the shared-autonomy belief estimator and alignment-based arbitration (Section~\ref{sec:sa}), the grasp lifecycle and world-model maintenance (Section~\ref{sec:grasp}), and the haptic interface and force rendering (Section~\ref{sec:teleop}). Chapter~\ref{ch:exp} validates the architecture on physical hardware (Section~\ref{sec:hwvalidation}), describes the setup of the user study (Section~\ref{sec:study}), and reports the results of the human-subject study. Chapter~\ref{ch:discussion} reads those results against the hypotheses of Section~\ref{sec:rqs}, weighs alternative explanations, and states threats to validity. Chapter~\ref{ch:conclusions} synthesises the whole, states the architecture's limitations, and outlines directions for future work; Appendix~\ref{app:config} collects every configuration parameter referenced throughout.
